\section{Performance Verbesserungen}

\subsection{Alpha-Beta-Pruning}
Zunächst war vorgesehen, sämtliche Optimierungen sowohl für die sequenzielle als auch für die parallele Implementierung des Minimax-Algorithmus umzusetzen.
Nach einer Recherche in Alpha-Beta-Pruning ergab sich, dass diese Optimierung für sequentielle Implementierungen vorgesehen ist.

Es bestehen auch Ansätze für parallele Implementierungen, wie Young Brothers Wait Concept.
Die Idee ist, dass der Zug mit der höchsten Wahrscheinlichkeit, der beste zu sein, zuerst evaluiert wird - danach werden alle anderen Züge in parallel evaluiert.
Somit können die restlichen Züge mit dem ersten verglichen werden\ \cite{alphabetaybwc}.

Bei solchen Ansätzen ist jedoch die Umsetzung umständlich und der Leistungsgewinn sehr variabel und unerheblich.
Aus diesen Gründen wurde die Alpha-Beta-Pruning Optimierung nur für die sequenzielle Implementierung umgesetzt.

Die ursprüngliche Implementierung des Minimax-Algorithmus musste nur minimal angepasst werden, um diese Optimierung einzubauen.
Es wurden 2 neue Parameter zur Funktion hinzugefügt: alpha und beta.
Diese merken sich die relevanten Werte der vorherigen Evaluationen und anhand dessen wird entschieden, ob die Suche vorzeitig abgebrochen werden kann (\cite{alphabetachessprogramming}).
Die relevanten Zeilen sind die Zeilen 66-75.
\lstinputlisting[
    firstline=30,
    lastline=78,
    firstnumber=30,
    caption={Minimax sequenziell mit Alpha-Beta-Pruning},
    label={lst:minimax-sequenziell-alpha-beta}
]{/C:/Users/romeo/IdeaProjects/blobfish/src/main/java/ch/hslu/cas/msed/blobfish/minimax/MiniMaxAlphaBetaSequential.java}

Zusätzlich ist bei dieser Optimierung die Reihenfolge der Züge wichtig.
Darum werden sie auf den Zeilen 45--46 sortiert.
Die Sortierung ist zunächst sehr primitiv - es werden Züge, die eine Figur fangen, zuerst ausgewertet, da diese eine grössere Chance haben bessere Werte zu liefern.
Mit dem Einbau weiterer Heuristiken wird auch diese Sortierung erweitert.